% Figure 1.2 --- RAG probabiliste contre accès déterministe via MCP
\begin{tikzpicture}[etape/.style={fbase, minimum height=1.5cm, text width=2.5cm}]

% ------------------------------------------------------------ RAG (axe 2)
\node[anchor=west, font=\bfseries\large, text=figbordeaux] at (0,1.55cm)
     {RAG (axe 2) --- probabiliste};

\foreach \i/\t in {0/{Requête de\\ l'analyste},
                   1/{Plongement\\ vectoriel},
                   2/{Base vecto-\\rielle : top-$k$\\ par similarité},
                   3/{$k$ passages\\ « voisins »},
                   4/{LLM condi-\\tionné},
                   5/{Réponse}}
  \node[etape, draw=figbordeaux, fill=figbordeaux!4, anchor=west]
       at ({\i*3.1cm},0) (r\i) {\t};

\foreach \i in {0,...,4} {\pgfmathtruncatemacro{\j}{\i+1}
  \draw[fgrosse] (r\i.east) -- (r\j.west);}

\node[below=4pt of r2.south, font=\small, text=figbordeaux, anchor=north]
     {surface d'attaque : la \emph{base} (5 textes $\Rightarrow$ 90 \%)};

% ------------------------------------------------------------ MCP (axe 3)
\node[anchor=west, font=\bfseries\large, text=figvert] at (0,-3.1cm)
     {Accès déterministe via MCP (axe 3)};

\foreach \i/\t in {0/{Requête de\\ l'analyste},
                   1/{LLM : décide\\ l'appel},
                   2/{Appel structuré\\ \texttt{get\_technique}\\ \texttt{("T1003")}},
                   3/{Serveur MCP :\\ index précons-\\truit, $O(1)$},
                   4/{Enregistrement\\ exact (STIX)},
                   5/{Réponse tracée}}
  \node[etape, draw=figvert, fill=figvert!5, anchor=west]
       at ({\i*3.1cm},-4.75cm) (m\i) {\t};

\foreach \i in {0,...,4} {\pgfmathtruncatemacro{\j}{\i+1}
  \draw[fgrosse] (m\i.east) -- (m\j.west);}

\node[below=4pt of m3.south, font=\small, text=figvert, anchor=north,
      align=center]
     {reproductible, journalisable ; sur-\\face d'attaque : le \emph{serveur} (axe 4)};

\end{tikzpicture}
